\section{The ordinary part of a maximal alternating extraction layer}
\label{sec:maximal-alternating-layer}

\paragraph{Accepted status and provenance.}
HD-ML01, question Q2dy, is a root synthesis of the accepted
D-TD11--13, HD-FT03/16--17/20--22, HD-TW04--07 and HD-AST13--15.
It introduces no new character estimate or finite computation. The
complete root reading and separate whole-proof review both pass
without mathematical repair. The arithmetic cutoff remains the already
effective HD-TW03 integer. No independent movement credit is claimed
for this synthesis.

Put
\[
 M_*=1+241\cdot240!\cdot2^{239},\qquad F_*=240M_*+2,
 \qquad N_*=[100(1+M_*)F_*^{240}]^2.
 \tag{ML.1}
\]
An \emph{extraction layer} means the whole irreducible tensor factor
removed at one nonabelian minimal-normal-subgroup step in D-TD12.
If that minimal normal subgroup is $A_n^t$, the layer consists of all
$t$ quasisimple component factors in that step, not one arbitrarily
selected component.

\begin{theorem}[HD-ML01]
Let an actual finite primitive coefficient representation $J$ on $V$
occur in the accepted rank-nine sixfold source interface. Fix one
complete recursive decomposition D-TD12. Suppose an alternating
extraction layer has natural degree $n\ge N_*$, maximal among the
alternating natural degrees of \emph{all} layers of this same
decomposition. Then every such maximal layer is one of the following:
\begin{enumerate}
\item a single ordinary standard factor, of dimension $n-1$;
\item a layer of genuine spin factors, with nontrivial central action
on the double cover of $A_n$.
\end{enumerate}
In particular all nonstandard ordinary layers, including ordinary
split constituents, and all ordinary standard layers with at least
two transitively permuted components are excluded. The conclusion
allows every external tensor factor, scalar phase and original
coefficient field. It is a necessary source restriction, not an
existence assertion or a bound on every remaining factor order.
\end{theorem}

\begin{proof}
The genuine spin alternative is retained, so it suffices to treat
an ordinary layer. Its component representation is faithful on
$A_n$. That group contains the elementary abelian group formed by
the Klein four groups on $\lfloor n/4\rfloor$ disjoint sets of
four letters. A faithful complex representation of an elementary
abelian group of rank $2\lfloor n/4\rfloor$ needs at least that
many characters to separate its elements, so its dimension is
at least $2\lfloor n/4\rfloor>6$. Hence $\dim V=e>6$.
The already checked intermediate-cover Euler and period argument
places this actual rank-nine coefficient source over the sphere,
with at least twelve genuine coefficient inertias and
\[
        \sum_i\operatorname{rank}(1-T_i)=2e+6.
 \tag{ML.2}
\]
This source interface is used as a necessary condition only.

Write the entire selected extraction layer as $W^{\otimes t}$ after
fixed compatible identifications, and let $U$ denote its entire
complement, so $V=W^{\otimes t}\otimes U$. The complement includes
both earlier and later pieces of the fixed recursion. D-TD12 proves
that all their full matrix algebras remain normalized or permuted
by the original image and that no multiplicity space is omitted.
The selected projective factor image is finite, contains its actual
minimal normal subgroup $A_n^t$, and permutes those components
transitively. These are properties of actual projected images;
the statement does not require the later components to lift to
independent normal subgroups of $J$.

Let $m$ be the largest alternating natural degree among the pieces
of $U$, with $m=1$ when no such piece occurs. By maximality,
$m\le n$. The full-compatible
normalizer and tensor-permutation bound HD-FT21 therefore gives
\[
 \frac{\nu(S)}{\dim U}\ge
 \frac{1}{\max(80,m-1)}\ge\frac1n
 \quad\text{for every nonscalar }S\text{ in the partner image}.
 \tag{ML.3}
\]
Here $n\ge80$. For a one-dimensional or trivial partner the
condition is vacuous. No new decomposition of $U$ with an
uncontrolled alternating degree is substituted: the estimate uses
the same complete decomposition throughout.

We make the actual quotient actions explicit. Conjugation on the
actual minimal normal subgroup gives a homomorphism from the
selected projected image to
$\operatorname{Aut}(A_n^t)=S_n\wr S_t$, since $n\ge8$.
Its image contains the inner subgroup $A_n^t$ and is transitive on
the components. Composing with the genuine quotient map from $J$
gives actual actions on $t$ factors and on $nt$ letters. Both are
transitive: the inner factors are transitive inside each block,
and the component action joins the blocks. Consequently the
quotient-cover Riemann--Hurwitz inequalities used by TW04--07 and
AST14 are applicable. Some genuine original marks may project
trivially; those theorems retain them through (ML.3) or full scalar
rank, and do not assume twelve nontrivial natural projections.

At $n\ge8$ the accepted covering-group classification leaves
$A_n$ and its double cover as the possible quasisimple groups
covering this simple quotient. The exact covering input is
Hoffman--Humphreys, Chapter4 opening, printed p.~31, as located
and used in HD-FT20.\cite{BI:HHM} An irreducible constituent with
nontrivial central action is of genuine spin type. We retain that
case here. Otherwise the constituent is ordinary. The ordinary
irreducible classification gives either an extendible restriction
of a Specht module, or one of the two constituents of a
self-conjugate Specht module; this is precisely the ordinary
restriction alternative used and located in HD-FT03. Component transitivity preserves
this distinction. In the split case, fixed odd relabelings may
be needed to identify all the factors; HD-TW07 proves that their
compatible cycle products are inner after those identifications.
In the extendible case, transpose partitions have the same
alternating restriction, so the row/column-tail notation is
unambiguous. A trivial ordinary factor is impossible because
the represented quasisimple factor is faithful and nontrivial.

For clarity the one integer in (ML.1) dominates every threshold
used next. We have $M_*\ge241$, hence $F_*\ge241$ and
$100(1+M_*)\ge241$. It follows directly that
\[
          N_*\ge241^{482}>3500000,
          \qquad N_*>1000000,\ 65536,\ 80.
 \tag{ML.4}
\]

First suppose $t=1$. The actual projective normalizer of an
ordinary constituent gives an actual $A_n$ or $S_n$ natural
quotient; in the split case the compatible quotient is $A_n$.
This follows either directly from the above conjugation action
and Schur's lemma or from the normalizer interfaces already
proved in FT03 and TW07. HD-FT03 excludes every split ordinary
constituent in this range, with arbitrary partners. HD-FT16--17
excludes every extendible ordinary constituent whose minimum
row/column tail is at least three at $n\ge241^{482}$.

For an extendible tail-two constituent the effective part of
HD-FT22 applies, using the actual stable quotient just verified
and the complete partner (ML.3). Its four genuine partner-only
entries imply
\[
                  4n^2\le61\max(80,m-1)\le61(n-1).
 \tag{ML.5}
\]
This is impossible for $n\ge10^6$. There are no uncovered
bounded pieces requiring an additional parameter $A$, because
D-TD12 has supplied the full quasisimple/Heisenberg decomposition.
Tail zero is the trivial representation and has been excluded.
The only remaining ordinary tail is one, the standard module.

Now suppose $t\ge2$. For extendible ordinary factors of minimum
row/column tail at least two, HD-TW05 uses (ML.3) and excludes the
entire layer at $n\ge N_*$. For split ordinary constituents,
HD-TW07 excludes it at $n\ge65536$ with the same partner floor.
For ordinary standard constituents, HD-AST14 excludes it at
$n\ge15$ with (ML.3), or equivalently HD-AST15 supplies the
maximal-layer consequence directly at $n\ge80$. The actual
factor and natural quotient actions used in each case have
already been verified above. Every ordinary possibility with
$t\ge2$ is therefore excluded. Together with the $t=1$ analysis,
this proves the theorem.
\end{proof}

\paragraph{Boundary and project interpretation.}
The single-standard-factor case with a nontrivial partner is not
the previously settled pure scalar-standard case. In particular
the exact identity
$\operatorname{Ind}_{J_1}^{J}(U|_{J_1})=U\oplus
(\operatorname{Std}_n\otimes U)$ retains the $U$-cohomology
summand. For deeper extraction layers, even the assertion that
$U(J_1)=U(J)$ needs a separate argument and is not used here.
The spin case is retained without an ordinary Specht estimate.
No assertion about bounded nonalternating rank, extraspecial
parameters, faithful inducing-image order, original induction
index, field descent, actual polarization or integral cycles
follows merely from (ML.1)--(ML.5).

This is a restriction on finite-monodromy coefficient local
systems in the curve-supplier approach. It changes no spreading,
integral numerator or positivity input. Its relation to condition A
(cote16, scan050), condition I (scan049 and cote44, Theorem4), and
the Cayley--Hamilton projector (scans048--049) is the already
recorded one: an actual detecting cycle must still be supplied
before those division and propagation arguments apply.
